\worksheettitle
  {Домашняя практика: ответы}
  {\modulemark{M06-H} Названия и диагонали}

\begin{teacherbox}[Критерии]
  Объяснение опирается на последовательный обход и соседство вершин, а не
  на порядок букв в алфавите или внешний вид отрезка.
\end{teacherbox}

\section{Ответы}

\begin{enumerate}
  \item Верны $ABCDE$, $CDEAB$, $AEDCB$. В $ACBED$ уже переход $A-C$
    проходит не по стороне.
  \item $AB$ --- сторона; $AD$ --- диагональ; $EA$ --- сторона.
  \item Из $K$: $KM$ и $KN$.
  \item Треугольник $BCD$ и четырёхугольник $BDEA$.
  \item $AB$ и $BA$ --- один отрезок и одна сторона. $AC$ соединяет
    несоседние вершины, поэтому является диагональю.
\end{enumerate}

\begin{resultbox}[Маршрут]
  При нарушении обхода или смешении стороны и диагонали назначьте M06-R.
  Единичный пропуск при верном способе требует только короткой практики.
  После устойчивого результата можно предложить M06\(\star\).
\end{resultbox}
